\subsection*{The Appointed Propers in Full}

\noindent The Latin below is the complete recurring formulary at marginal
nos.~1582--1591 in the 1962 Vatican typical Missal, printed pp.~396--397.
Scriptural citations use the Vulgate numbering printed in the Missal. Three
elements materially adapt their biblical source forms.

\fulltextheading{Introit\enspace\properrefs{Int.}}\label{introit}
\fulltextcueelement{\latin{Psalmus 85:1 et 2--3; 4}; no.~1582}
{Inclína, Dómine, aurem tuam ad me, et exáudi me: salvum fac servum tuum,
Deus meus, sperántem in te: miserére mihi, Dómine, quóniam ad te clamávi
tota die.}{Ps. ibid., 4}{Lætífica ánimam servi tui: quia ad te,
Dómine, ánimam meam levávi. ℣.~Glória Patri.}
\fulltextgap{The antiphon joins selected clauses from vv.~1--3 and omits the
descriptions of the petitioner as needy, poor, and holy. The psalm verse is
v.~4.}

\fulltextheading{Collect\enspace\properrefs{Coll.}}\label{collect}
\fulltextelement{\latin{Oratio}; no.~1583}
Ecclésiam tuam, Dómine, miserátio continuáta mundet et múniat: et quia sine
te non potest salva consístere; tuo semper múnere gubernétur. Per Dóminum.
\finishfulltextelement

\Needspace{24\baselineskip}
\fulltextheading{Epistle\enspace\properrefs{Ep.}}\label{epistle}
\fulltextelement{\latin{Ad Galatas 5:25--26; 6:1--10}; no.~1584}
Léctio Epístolæ beáti Pauli Apóstoli ad Gálatas.

Fratres: Si spíritu vívimus, spíritu et ambulémus. Non efficiámur inánis
glóriæ cúpidi, ínvicem provocántes, ínvicem invidéntes. Fratres, et si
præoccupátus fúerit homo in áliquo delícto, vos, qui spirituáles estis,
huiúsmodi instrúite in spíritu lenitátis, consíderans teípsum, ne et tu
tentéris. Alter altérius ónera portáte, et sic adimplébitis legem Christi.
Nam si quis exístimat se áliquid esse, cum nihil sit, ipse se sedúcit. Opus
autem suum probet unusquísque, et sic in semetípso tantum glóriam habébit,
et non in áltero. Unusquísque enim onus suum portábit. Commúnicet autem is
qui catechizátur verbo, ei qui se catechízat, in ómnibus bonis. Nolíte
erráre: Deus non irridétur. Quæ enim semináverit homo, hæc et metet. Quóniam
qui séminat in carne sua, de carne et metet corruptiónem: qui autem séminat
in spíritu, de spíritu metet vitam ætérnam. Bonum autem faciéntes, non
deficiámus: témpore enim suo metémus, non deficiéntes. Ergo dum tempus
habémus, operémur bonum ad omnes, máxime autem ad domésticos fídei.
\finishfulltextelement

\fulltextheading{Gradual\enspace\properrefs{Grad.}}\label{gradual}
\fulltextelement{\latin{Psalmus 91:2--3}; no.~1585}
Bonum est confitéri Dómino: et psállere nómini tuo, Altíssime. ℣.~Ad
annuntiándum mane misericórdiam tuam et veritátem tuam per noctem.
\finishfulltextelement

\fulltextheading{Alleluia\enspace\properrefs{All.}}\label{alleluia}
\fulltextelement{\latin{Psalmus 94:3}; no.~1586}
Allelúia, allelúia. ℣.~Quóniam Deus magnus Dóminus, et Rex magnus super
omnem terram. Allelúia.
\finishfulltextelement
\fulltextgap{The Missal ends with \latin{super omnem terram}. The Clementine
Vulgate reads \latin{super omnes deos}; Challoner renders that ending
``above all gods.''}

\Needspace{25\baselineskip}
\fulltextheading{Gospel\enspace\properrefs{Gosp.}}\label{gospel}
\fulltextelement{\latin{Lucas 7:11--16}; no.~1587}
✠ Sequéntia sancti Evangélii secúndum Lucam.

In illo témpore: Ibat Iesus in civitátem, quæ vocátur Naim: et ibant cum eo
discípuli eius, et turba copiósa. Cum autem appropinquáret portæ civitátis,
ecce defúnctus efferebátur fílius únicus matris suæ: et hæc vídua erat: et
turba civitátis multa cum illa. Quam cum vidísset Dóminus, misericórdia motus
super eam, dixit illi: Noli flere. Et accéssit, et tétigit lóculum. (Hi
autem, qui portábant, stetérunt). Et ait: Aduléscens, tibi dico, surge. Et
resédit qui erat mórtuus, et cœpit loqui. Et dedit illum matri suæ. Accépit
autem omnes timor: et magnificábant Deum, dicéntes: Quia prophéta magnus
surréxit in nobis: et quia Deus visitávit plebem suam.
\finishfulltextelement
\par\smallskip\latin{Credo.}
\fulltextgap{The Missal supplies the liturgical opening \latin{In illo
témpore}; Luke's Vulgate text begins \latin{Et factum est: deinceps}. The
Creed belongs to the Ordinary.}

\fulltextheading{Offertory\enspace\properrefs{Off.}}\label{offertory}
\fulltextelement{\latin{Psalmus 39:2, 3 et 4}; no.~1588}
Exspéctans exspectávi Dóminum, et respéxit me: et exaudívit deprecatiónem
meam: et immísit in os meum cánticum novum, hymnum Deo nostro.
\finishfulltextelement
\fulltextgap{The antiphon joins clauses from vv.~2--4 and substitutes
\latin{respéxit me}, \latin{deprecatiónem meam}, and \latin{hymnum} for the
Clementine forms. The pit, mire, rock, steadied steps, and many who see are
complete-context material rather than appointed words.}

\fulltextheading{Secret\enspace\properrefs{Sec.}}\label{secret}
\fulltextelement{\latin{Secreta; Præfatio de Sanctissima Trinitate}; no.~1589}
Tua nos, Dómine, sacraménta custódiant: et contra diabólicos semper tueántur
incúrsus. Per Dóminum nostrum Iesum Christum, Fílium tuum: Qui tecum vivit
et regnat in unitáte.

Præfatio de Ssma Trinitate.
\finishfulltextelement

\fulltextheading{Communion\enspace\properrefs{Comm.}}\label{communion}
\fulltextelement{\latin{Ioannes 6:52}; no.~1590}
Panis, quem ego dédero, caro mea est pro sæculi vita.
\finishfulltextelement
\fulltextgap{The antiphon appoints only the final bread-and-flesh clause. It
reads \latin{pro sæculi vita}; the Clementine verse has \latin{pro mundi
vita}. Vulgate John~6:52 corresponds to modern 6:51b.}

\Needspace{12\baselineskip}
\fulltextheading{Postcommunion\enspace\properrefs{Postcomm.}}\label{postcommunion}
\fulltextelement{\latin{Postcommunio}; no.~1591}
Mentes nostras et córpora possídeat, quǽsumus, Dómine, doni cæléstis
operátio: ut non noster sensus in nobis, sed iúgiter eius prævéniat efféctus.
Per Dóminum nostrum Iesum Christum, Fílium tuum: Qui tecum vivit et regnat in
unitáte.
\finishfulltextelement
